\begin{abstract}
Indoor environments present significant safety hazards including
slippery floor surfaces, furniture and clutter obstacles, and sharp
objects, which collectively represent a leading cause of domestic
injuries worldwide. While semantic segmentation provides a natural
framework for pixel-level hazard localisation, existing architectures
predict all semantic classes simultaneously from a single decoder,
providing no mechanism to exploit the hierarchical structure inherent
to hazard-aware scene understanding.

This paper presents the \textbf{Dual-Stream Progressive Hazard Network
(DSPHN)}, whose primary objective is \textit{architectural innovation
for hazard-aware scene understanding} rather than solely maximising
segmentation accuracy. DSPHN introduces three architectural
contributions: (1)~a \textbf{Dual-Stream Appearance-Context Decoder}
that disentangles local visual hazard cues from spatial contextual
relationships using parallel standard and dilated convolutional
decoder pathways; (2)~\textbf{Severity-Guided Adaptive Fusion Gates}
that blend the two decoder streams using gates initialised from
domain hazard severity scores; and (3)~a \textbf{Two-Stage Hazard
Refinement Mechanism} where Stage~1 produces a binary safe/unsafe
map that serves as a direct structural input to Stage~2 fine
classification.

To the best of our knowledge, this combination of dual-stream
appearance-context disentanglement with hierarchical safety-aware
conditional inference has not been previously proposed for indoor
hazard segmentation. We further contribute a four-class indoor hazard
taxonomy and 100~manually annotated images. Experiments on a
3,000-image hybrid dataset yield test Mean~IoU of \textbf{0.4198}
and Pixel~Accuracy of \textbf{0.8332}. A controlled ablation study
confirms that the full DSPHN outperforms both the dual-stream
component alone (mIoU~0.3894) and the two-stage component alone
(mIoU~0.3793), validating the contribution of each architectural
innovation.
\end{abstract}
